\begin{textsample}{NEG Dim 2 – vem\_ed\_04 – Score 6.00 – vem\_ed\_04/t016.txt}  \label{ex:f2_neg_133}
BOAS PRÁTICAS Com o pé direito Fatec Ipiranga estreia nos intercâmbios virtuais com um projeto desenvolvido com a Jamestown Community College, ligada à State University of New York (JCC / SUNY).
As atividades foram conduzidas por José Carlos Barbosa Lopes, que leciona Inglês na Fatec Ipiranga, e Renee Funke, da área de Educação da JCC / SUNY.
No primeiro encontro, em 18 de setembro, os 26 alunos brasileiros perguntaram via Zoom sobre a experiência profissional da professora norte-americana, a adaptação às aulas online e a empregabilidade na pandemia.
Para preparar os estudantes, Lopes oferece oficinas de conversação em inglês – além das duas horas-aula semanais do idioma obrigatórias nos currículos das Fatecs, diferencial ante outras instituições de ensino superior tecnológico.
Segundo Lopes, “o projeto é importante para os alunos da Fatec Ipiranga, não só pelo aprimoramento da língua inglesa, mas também pelo contato com estudantes de outro país e a experiência de discutir assuntos de interesse global”.
Sobre o planejamento do intercâmbio virtual, o professor afirma que foram considerados os perfis do alunado das duas instituições, os objetivos das disciplinas e as temáticas afins para compor a colaboração.
Na fase “quebra-gelo”, os 40 participantes (no Brasil e nos EUA) elaboraram um vídeo pessoal, trocaram mensagens e definiram os assuntos mais relevantes para o grupo – e assim foram montadas as equipes mistas (Fatec Ipiranga e JCC / SUNY).
As ferramentas digitais usadas foram Zoom, Padlet, Facebook, WhatsApp e Teams.
Como a maioria dos 26 participantes da Fatec Ipiranga frequenta aulas e oficinas de Inglês, foi possível “discutir e avaliar nos grupos a importância de fatores como colaboração, relações interpessoais, aspectos culturais e recursos de negociação de sentido, todos fundamentais num projeto de intercâmbio internacional”, relata Lopes. “O fato de poderem participar de discussões em inglês sobre assuntos tão relevantes atualmente como pandemia, multiculturalidade, questões sociais e políticas, foi inspirador”.
Perguntado sobre que dica daria para professores que desejam começar um intercâmbio virtual, Lopes recomenda: “Creio que seja interessante observar as possibilidades de aperfeiçoamento que poderão ser agregadas ao curso, à formação dos alunos e à disciplina”.
Cabe lembrar que a equipe dos PCIs pode ajudar em todos esses passos.

% matched lemmas: 
\end{textsample}
